% $Log: abstract.tex,v $
% Revision 1.1  93/05/14  14:56:25  starflt
% Initial revision
% 
% Revision 1.1  90/05/04  10:41:01  lwvanels
% Initial revision
% 
%
%% The text of your abstract and nothing else (other than comments) goes here.
%% It will be single-spaced and the rest of the text that is supposed to go on
%% the abstract page will be generated by the abstractpage environment.  This
%% file should be \input (not \include 'd) from cover.tex.

The fast development of Unmanned Aerial Vehicles (UAVs) in civil and military applications demands control systems that ensure high reliability and safety, particularly in the presence of system uncertainties, external disturbances, and physical faults. This thesis addresses the design of advanced fault-tolerant control (FTC) strategies for various autonomous aerial platforms, ranging from grounded experimental prototypes to complex multi-agent swarms. The core objective is to guarantee trajectory tracking with fixed-time convergence, ensuring that system errors stabilize within a pre-calculated time independent of initial conditions.

The research is divided into three main contributions. First, a fault-tolerant fixed-time adaptive neural controller is proposed for a Twin Rotor MIMO System (TRMS). This approach utilizes Radial Basis Function Neural Networks (RBFNN) to approximate unknown nonlinear dynamics and sensor faults, with control parameters optimized via the Grey Wolf Optimizer (GWO). Second, the work extends to a quadrotor UAV, developing an adaptive fuzzy backstepping control scheme. Here, Fuzzy Logic Systems (FLS) are employed to compensate for actuator faults and aerodynamic disturbances, while the Whale Optimization Algorithm (WOA) is utilized to tune the controller gains for optimal performance. Finally, the thesis tackles the challenge of Coaxial UAV Swarms, introducing a distributed control protocol for formation flight. A novel Self-Structuring Neural Network (SSNN) is integrated to dynamically handle structural uncertainties and faults within the swarm, allowing for the autonomous addition or pruning of neurons based on system complexity. Extensive numerical simulations demonstrate the effectiveness, robustness, and superior tracking accuracy of the proposed strategies compared to conventional methods.


\textbf{Keywords:} 
Fault-Tolerant Control (FTC), Unmanned Aerial Vehicles (UAV), Fixed-Time Stability, Neural Networks, Fuzzy Logic, Swarm Formation, Optimization Algorithms (GWO, WOA).



\setcode{utf8}
\vspace{6pt}
\begin{RLtext}\small

\textbf{الملخص}

إن الانتشار السريع للمركبات الجوية غير المقادة  في التطبيقات المدنية والعسكرية يتطلب أنظمة تحكم تضمن مستويات عالية من الموثوقية والأمان، خاصة في ظل وجود عدم اليقين في النظام، الاضطرابات الخارجية، والأعطال الفيزيائية. تتناول هذه الأطروحة تصميم استراتيجيات متقدمة لالتحكم المتسامح مع الأخطاء  لمختلف المنصات الجوية المستقلة، بدءاً من النماذج التجريبية المثبتة وصولاً إلى أسراب الطائرات المعقدة. الهدف الرئيسي هو ضمان تتبع المسار مع التقارب في زمن ثابت ، مما يضمن استقرار أخطاء النظام خلال وقت محسوب مسبقاً ومستقل عن الشروط الابتدائية.

ينقسم البحث إلى ثلاث مساهمات رئيسية. أولاً، تم اقتراح متحكم عصبي تكيفي ذو زمن ثابت ومقاوم للأعطال لنظام الدوّار المزدوج . تعتمد هذه المقاربة على الشبكات العصبية ذات الدالة القطرية  لتقدير الديناميكيات غير الخطية المجهولة وأعطال الحساسات، مع تحسين بارامترات التحكم باستخدام خوارزمية الذئب الرمادي . ثانياً، توسع العمل ليشمل طائرة كوادروت ، حيث تم تطوير نظام تحكم تراجعي  يعتمد على المنطق الضبابي. تم استخدام أنظمة المنطق الضبابي  لتعويض أعطال المحركات والاضطرابات الديناميكية الهوائية، بينما استُخدمت خوارزمية تحسين الحوت  لضبط مكاسب المتحكم. أخيراً، تعالج الأطروحة تحدي أسراب الطائرات المحورية ، مقدمةً بروتوكول تحكم موزع للطيران التشكيلي. تم دمج شبكة عصبية ذاتية الهيكلة  للتعامل ديناميكياً مع الشكوك الهيكلية والأعطال داخل السرب. أثبتت المحاكاة الرقمية المكثفة فعالية ومتانة ودقة الاستراتيجيات المقترحة مقارنة بالطرق التقليدية.

\textbf{الكلمات المفتاحية}
الكلمات المفتاحية: التحكم المتسامح مع الأخطاء، الطائرات بدون طيار ، الاستقرار في زمن ثابت، الشبكات العصبية، المنطق الضبابي، تشكيل الأسراب، خوارزميات التحسين .

\end{RLtext}
